% !TEX spellcheck = en_US



% ~~~~~~~~~~~~~~~~~~~~~~~~~~~~~~~~~~~~~~~ V
% (NC) 2026 Haluk Bingol
% github.com/halukbingol/LaTeX-Templates
%
% Licensees may copy, distribute, display, and perform the work and make derivative works 
% and remixes based on it only for non-commercial purposes. 
% ~~~~~~~~~~~~~~~~~~~~~~~~~~~~~~~~~~~~~~~ A




% =====================================================================
%  Trees -- The Concepts
%  ---------------------------------------------------------------------
%  Built on hblib/Template/hblibTemplate.tex with the libraries in ../../hblib/.
%  Programs and their real outputs are in codeoutput/:
%     codeoutput/<Name>.java  a program (classes without main() are only compiled)
%     codeoutput/<Name>.sh    a terminal session, run in a copy of codeoutput/<Name>/
%     codeoutput/<Name>.txt   what it printed, made by:  bash run_examples.sh
% =====================================================================

\documentclass[aspectratio=169,10pt,compress]{beamer}

% ~~~~~~~~~~~~~~~~~~~~~~~~~~~~~~~~~~~~~~~~~~~~~~~~~~~~~~~~~~~~~~~~~ V hblib
	\usepackage{../../hblib/hblibPresentation} % lib presentation: theme, i/N, \hSection, algorithm2e, ...
	\usepackage{../../hblib/hblibCodeoutput} % lib code + output: \lstcode, \lstterminal, ...
	\usepackage{../../hblib/hblibDatastructures} % lib data structures: \hString, \hStack, \hQueue
% ~~~~~~~~~~~~~~~~~~~~~~~~~~~~~~~~~~~~~~~~~~~~~~~~~~~~~~~~~~~~~~~~~ A hblib


% xcolor, array (and the L{width} column), listings, tikz and algorithm2e come from the hblib libraries
\usepackage[T1]{fontenc}
\usepackage{lmodern}
\usepackage{booktabs}

\newcommand{\win}{\textbf{Windows}}
\newcommand{\mac}{\textbf{macOS}}

% algorithm2e set-up (the package is loaded by hblibPresentation)
\SetAlCapFnt{\small}
\SetAlCapNameFnt{\small}
\SetKwProg{Fn}{function}{}{end}
\SetKw{KwError}{error}

% ---------------------------------------------------------------------
\title{Trees}
\subtitle{The Concepts}
\author[Bingol]{Haluk O. Bingol}

\institute[FCIS]{
    Faculty of Computer and Information Sciences\\
    Yeditepe University
}
\date{
	{\tiny \hbTimeStamp}
}

\begin{document}




% ~~~~~~~~~~~~~~~~~~~~~~~~~~~~~~~~~~~~~~ V frm
\begin{frame}[t,fragile]
  \titlepage
\end{frame}
% ~~~~~~~~~~~~~~~~~~~~~~~~~~~~~~~~~~~~~~ A frm




% ~~~~~~~~~~~~~~~~~~~~~~~~~~~~~~~~~~~~~~ V frm
\begin{frame}[t,fragile] \frametitle{Outline}
  \begin{columns}[T]
    \begin{column}{0.48\textwidth}
      \tableofcontents[sections={1-5}, hideallsubsections]
    \end{column}
    \begin{column}{0.48\textwidth}
      \tableofcontents[sections={6-9}, hideallsubsections]
    \end{column}
  \end{columns}
\end{frame}
% ~~~~~~~~~~~~~~~~~~~~~~~~~~~~~~~~~~~~~~ A frm

% ~~~~~~~~~~~~~~~~~~~~~~~~~~~~~~~~~~~~~~ V
\hSection[Trees]{What Is a Tree?}
% ~~~~~~~~~~~~~~~~~~~~~~~~~~~~~~~~~~~~~~ A




% ~~~~~~~~~~~~~~~~~~~~~~~~~~~~~~~~~~~~~~ V frm
\begin{frame}[t,fragile] \frametitle{Trees are everywhere}
  \begin{columns}[T]
    \begin{column}{0.5\textwidth}
      \begin{center}
      \begin{tikzpicture}[htree, level distance=0.95cm, level 1/.style={sibling distance=2.2cm}, level 2/.style={sibling distance=1.1cm}, scale=0.85, transform shape]
  \node[hsquare] {/} child {node[hsquare] {Users} child {node[hsquare] {ali}} child {node[hsquare] {ayse}}} child {node[hsquare] {Applications}} child {node[hsquare] {System} child {node[hsquare] {Library}}};
\end{tikzpicture}

      \small a file system (\mac: \texttt{/Users/ali}; \win: \texttt{C:\textbackslash Users\textbackslash ali})
      \end{center}
    \end{column}
    \begin{column}{0.46\textwidth}
      \begin{itemize}\small
        \item folders inside folders
        \item a company: managers and their teams
        \item a web page: \texttt{<html>} contains \texttt{<body>} contains \texttt{<p>} \dots{} (the DOM)
        \item an arithmetic expression: \texttt{(3 + 4) * (5 - 2)}
        \item a family tree, a tournament, a book with chapters and sections
      \end{itemize}
      \vspace{0.3em}
      \small Data where each item belongs to \textbf{exactly one} ``parent'' item: a \textbf{hierarchy}.
    \end{column}
  \end{columns}
\end{frame}
% ~~~~~~~~~~~~~~~~~~~~~~~~~~~~~~~~~~~~~~ A frm




% ~~~~~~~~~~~~~~~~~~~~~~~~~~~~~~~~~~~~~~ V frm
\begin{frame}[t,fragile] \frametitle{Definition}
  \begin{columns}[T]
    \begin{column}{0.52\textwidth}
      \begin{block}{Recursive definition}
        A \textbf{tree} is either
        \begin{itemize}
          \item \textbf{empty}, or
          \item a \textbf{root} node with zero or more \textbf{subtrees}, each of them a tree.
        \end{itemize}
      \end{block}
      \begin{block}{As a graph}
        \small A connected graph with no cycles; one node is chosen as the root. Every node except the root has exactly one parent.
      \end{block}
    \end{column}
    \begin{column}{0.44\textwidth}
      \begin{center}
      \begin{tikzpicture}[htree, level distance=0.8cm, level 1/.style={sibling distance=2.2cm}, level 2/.style={sibling distance=1.1cm}, level 3/.style={sibling distance=0.6cm}, scale=0.75, transform shape]
  \node[hnode] {A} child {node[hnode] {B} child {node[hnode] {D}} child {node[hnode] {E}}} child {node[hnode] {C} child[missing] child {node[hnode] {F} child {node[hnode] {G}} child[missing]}};
\end{tikzpicture}\\
      \small \textbf{tree}: 7 nodes, 6 edges
      \end{center}
      \begin{center}
      \begin{tikzpicture}[scale=0.7, transform shape]
        \node[hnode] (x) at (0,0) {X}; \node[hnode] (y) at (-0.8,-0.9) {Y}; \node[hnode] (z) at (0.8,-0.9) {Z};
        \node[hnode] (w) at (0,-1.8) {W};
        \draw[thick, gray!70] (x) -- (y) (x) -- (z) (y) -- (w) (z) -- (w);
      \end{tikzpicture}\\
      \small \hRemark{not a tree}: \texttt{W} has two parents (a cycle)
      \end{center}
    \end{column}
  \end{columns}
\end{frame}
% ~~~~~~~~~~~~~~~~~~~~~~~~~~~~~~~~~~~~~~ A frm

% ~~~~~~~~~~~~~~~~~~~~~~~~~~~~~~~~~~~~~~ V
\hSection[Terms]{Terminology}
% ~~~~~~~~~~~~~~~~~~~~~~~~~~~~~~~~~~~~~~ A




% ~~~~~~~~~~~~~~~~~~~~~~~~~~~~~~~~~~~~~~ V frm
\begin{frame}[t,fragile] \frametitle{Words for the parts of a tree}
  \begin{columns}[T]
    \begin{column}{0.4\textwidth}
      \begin{center}
      \begin{tikzpicture}[htree, level distance=0.95cm, level 1/.style={sibling distance=2.2cm}, level 2/.style={sibling distance=1.1cm}, level 3/.style={sibling distance=0.6cm}, scale=1.0, transform shape]
  \node[hnode, hhot] {A} child {node[hnode] {B} child {node[hnode, hleaf] {D}} child {node[hnode, hleaf] {E}}} child {node[hnode] {C} child[missing] child {node[hnode] {F} child {node[hnode, hleaf] {G}} child[missing]}};
\end{tikzpicture}
      \end{center}
      \small green: the root; yellow: the leaves
    \end{column}
    \begin{column}{0.58\textwidth}
      \small
      \begin{tabular}{lL{0.62\textwidth}}
        \toprule
        \textbf{root}            & the top node, no parent: \texttt{A} \\
        \textbf{parent, child}   & \texttt{B} is the parent of \texttt{D}; \texttt{D} is a child of \texttt{B} \\
        \textbf{siblings}        & same parent: \texttt{D} and \texttt{E} \\
        \textbf{leaf}            & no children: \texttt{D}, \texttt{E}, \texttt{G} \\
        \textbf{internal node}   & has a child: \texttt{A}, \texttt{B}, \texttt{C}, \texttt{F} \\
        \textbf{ancestors}       & on the way up to the root: of \texttt{G} are \texttt{F}, \texttt{C}, \texttt{A} \\
        \textbf{descendants}     & everything below: of \texttt{C} are \texttt{F}, \texttt{G} \\
        \textbf{subtree}         & a node with all its descendants \\
        \textbf{degree}          & number of children: \texttt{deg(B) = 2}, \texttt{deg(C) = 1} \\
        \bottomrule
      \end{tabular}
    \end{column}
  \end{columns}
\end{frame}
% ~~~~~~~~~~~~~~~~~~~~~~~~~~~~~~~~~~~~~~ A frm




% ~~~~~~~~~~~~~~~~~~~~~~~~~~~~~~~~~~~~~~ V frm
\begin{frame}[t,fragile] \frametitle{Depth, level and height}
  \begin{columns}[T]
    \begin{column}{0.4\textwidth}
      \begin{center}
      \begin{tikzpicture}[htree, level distance=0.95cm, level 1/.style={sibling distance=2.2cm}, level 2/.style={sibling distance=1.1cm}, level 3/.style={sibling distance=0.6cm}, scale=1.0, transform shape]
  \node[hnode] {A} child {node[hnode] {B} child {node[hnode] {D}} child {node[hnode] {E}}} child {node[hnode] {C} child[missing] child {node[hnode] {F} child {node[hnode] {G}} child[missing]}};
\end{tikzpicture}
      \end{center}
    \end{column}
    \begin{column}{0.58\textwidth}
      \begin{itemize}\small
        \item \textbf{depth} of a node: number of edges from the root. \texttt{depth(A) = 0}, \texttt{depth(G) = 3}.
        \item \textbf{level} $k$: all nodes with depth $k$. Level 1: \texttt{B}, \texttt{C}.
        \item \textbf{height} of a node: edges on the longest path down to a leaf. \texttt{height(C) = 2}, \texttt{height(D) = 0}.
        \item \textbf{height of the tree} = height of the root = 3.
        \item The empty tree has height $-1$ (so that a single node has height 0).
      \end{itemize}
      \begin{center}\small
      \begin{tabular}{lccccccc}
        \toprule
        node   & A & B & C & D & E & F & G \\
        \midrule
        depth  & 0 & 1 & 1 & 2 & 2 & 2 & 3 \\
        height & 3 & 1 & 2 & 0 & 0 & 1 & 0 \\
        \bottomrule
      \end{tabular}
      \end{center}
    \end{column}
  \end{columns}
\end{frame}
% ~~~~~~~~~~~~~~~~~~~~~~~~~~~~~~~~~~~~~~ A frm




% ~~~~~~~~~~~~~~~~~~~~~~~~~~~~~~~~~~~~~~ V frm
\begin{frame}[t,fragile] \frametitle{Two facts}
  \begin{block}{A tree with $n$ nodes has $n - 1$ edges}
    \small Every node except the root has exactly one edge to its parent.
  \end{block}
  \begin{block}{Between any two nodes there is exactly one path}
    \small Go up from both to their lowest common ancestor, then down. Two different paths would make a cycle.
  \end{block}
  \vspace{0.3em}
  \small Height and size are defined \textbf{recursively}, like the tree itself:
  \[
  	\mathrm{size}(t) = 1 + \sum_{c \text{ child of } t} \mathrm{size}(c)
  \]
  \[
  	\mathrm{height}(t) = 1 + \max_{c \text{ child of } t} \mathrm{height}(c), \qquad \mathrm{height}(\text{empty}) = -1
  \]
\end{frame}
% ~~~~~~~~~~~~~~~~~~~~~~~~~~~~~~~~~~~~~~ A frm

% ~~~~~~~~~~~~~~~~~~~~~~~~~~~~~~~~~~~~~~ V
\hSection[Binary]{Binary Trees}
% ~~~~~~~~~~~~~~~~~~~~~~~~~~~~~~~~~~~~~~ A




% ~~~~~~~~~~~~~~~~~~~~~~~~~~~~~~~~~~~~~~ V frm
\begin{frame}[t,fragile] \frametitle{Binary tree: at most two children, left and right}
  \begin{columns}[T]
    \begin{column}{0.5\textwidth}
      \begin{block}{Definition}
        A \textbf{binary tree} is either empty, or a root with a \textbf{left} subtree and a \textbf{right} subtree, both binary trees.
      \end{block}
      \begin{itemize}\small
        \item Left and right are \textbf{different}: the two trees on the right are not the same binary tree.
        \item Each node has 0, 1 or 2 children.
      \end{itemize}
    \end{column}
    \begin{column}{0.46\textwidth}
      \begin{center}
      \begin{tikzpicture}[htree, level distance=0.95cm, level 1/.style={sibling distance=1.4cm}, scale=1.0, transform shape]
  \node[hnode] {P} child {node[hnode] {Q}} child[missing];
\end{tikzpicture}\qquad
      \begin{tikzpicture}[htree, level distance=0.95cm, level 1/.style={sibling distance=1.4cm}, scale=1.0, transform shape]
  \node[hnode] {P} child[missing] child {node[hnode] {Q}};
\end{tikzpicture}\\[0.3em]
      \small \texttt{Q} is the \textbf{left} child \quad\quad \texttt{Q} is the \textbf{right} child
      \end{center}
    \end{column}
  \end{columns}
\end{frame}
% ~~~~~~~~~~~~~~~~~~~~~~~~~~~~~~~~~~~~~~ A frm




% ~~~~~~~~~~~~~~~~~~~~~~~~~~~~~~~~~~~~~~ V frm
\begin{frame}[t,fragile] \frametitle{Shapes of binary trees}
  \begin{columns}[T]
    \begin{column}{0.24\textwidth}
      \centering
      \begin{tikzpicture}[htree, level distance=0.95cm, level 1/.style={sibling distance=1.2cm}, level 2/.style={sibling distance=0.6cm}, scale=0.75, transform shape]
  \node[hnode] {1} child {node[hnode] {2} child {node[hnode] {4}} child {node[hnode] {5}}} child {node[hnode] {3}};
\end{tikzpicture}\\
      \small \textbf{full}: every node has 0 or 2 children
    \end{column}
    \begin{column}{0.24\textwidth}
      \centering
      \begin{tikzpicture}[htree, level distance=0.95cm, level 1/.style={sibling distance=1.2cm}, level 2/.style={sibling distance=0.6cm}, scale=0.75, transform shape]
  \node[hnode] {1} child {node[hnode] {2} child {node[hnode] {4}} child {node[hnode] {5}}} child {node[hnode] {3} child {node[hnode] {6}} child[missing]};
\end{tikzpicture}\\
      \small \textbf{complete}: all levels full, the last filled from the left
    \end{column}
    \begin{column}{0.24\textwidth}
      \centering
      \begin{tikzpicture}[htree, level distance=0.95cm, level 1/.style={sibling distance=1.2cm}, level 2/.style={sibling distance=0.6cm}, scale=0.75, transform shape]
  \node[hnode] {1} child {node[hnode] {2} child {node[hnode] {4}} child {node[hnode] {5}}} child {node[hnode] {3} child {node[hnode] {6}} child {node[hnode] {7}}};
\end{tikzpicture}\\
      \small \textbf{perfect}: all leaves at the same depth
    \end{column}
    \begin{column}{0.24\textwidth}
      \centering
      \begin{tikzpicture}[htree, level distance=0.6cm, level 1/.style={sibling distance=0.8cm}, level 2/.style={sibling distance=0.8cm}, level 3/.style={sibling distance=0.8cm}, scale=0.75, transform shape]
  \node[hnode] {1} child[missing] child {node[hnode] {2} child[missing] child {node[hnode] {3} child[missing] child {node[hnode] {4}}}};
\end{tikzpicture}\\
      \small \textbf{degenerate}: a list in disguise
    \end{column}
  \end{columns}
  \vspace{0.5em}
  \small (Some books use ``complete'' for what is called ``perfect'' here; check the definitions in your textbook.)
\end{frame}
% ~~~~~~~~~~~~~~~~~~~~~~~~~~~~~~~~~~~~~~ A frm




% ~~~~~~~~~~~~~~~~~~~~~~~~~~~~~~~~~~~~~~ V frm
\begin{frame}[t,fragile] \frametitle{How many nodes? How high?}
  For a binary tree of height $h$ with $n$ nodes:
  \[
  	\text{at most } 2^{k} \text{ nodes at level } k
  \]
  \[
  	h + 1 \le n \le 2^{h+1} - 1
  \]
  \[
  	\lfloor \log_{2} n \rfloor \le h \le n - 1
  \]
  \begin{itemize}\small
    \item A perfect tree of height 19 has $2^{20} - 1 \approx$ one million nodes: any node is at most 19 steps from the root.
    \item A degenerate tree with a million nodes has height 999\,999.
    \item Fast tree algorithms keep $h$ close to $\log_{2} n$: \textbf{balanced} trees (AVL, red-black; junior deck).
  \end{itemize}
\end{frame}
% ~~~~~~~~~~~~~~~~~~~~~~~~~~~~~~~~~~~~~~ A frm

% ~~~~~~~~~~~~~~~~~~~~~~~~~~~~~~~~~~~~~~ V
\hSection[N-ary]{N-ary Trees}
% ~~~~~~~~~~~~~~~~~~~~~~~~~~~~~~~~~~~~~~ A




% ~~~~~~~~~~~~~~~~~~~~~~~~~~~~~~~~~~~~~~ V frm
\begin{frame}[t,fragile] \frametitle{Trees with more children}
  \begin{columns}[T]
    \begin{column}{0.48\textwidth}
      \begin{center}
      \begin{tikzpicture}[htree, level distance=0.95cm, level 1/.style={sibling distance=2.0cm}, level 2/.style={sibling distance=0.75cm}, scale=1.0, transform shape]
  \node[hnode] {R} child {node[hnode] {A} child {node[hnode] {D}} child {node[hnode] {E}}} child {node[hnode] {B}} child {node[hnode] {C} child {node[hnode] {F}} child {node[hnode] {G}} child {node[hnode] {H}}};
\end{tikzpicture}
      \end{center}
      \small a general tree: \texttt{C} has 3 children, \texttt{B} none
    \end{column}
    \begin{column}{0.48\textwidth}
      \begin{itemize}\small
        \item \textbf{$k$-ary tree}: every node has at most $k$ children (binary: $k = 2$, ternary: $k = 3$).
        \item \textbf{general} (N-ary) tree: any number of children, kept in an \textbf{ordered list}.
        \item At most $k^{d}$ nodes at depth $d$; a perfect $k$-ary tree of height $h$ has $\frac{k^{h+1} - 1}{k - 1}$ nodes.
        \item Examples: folders, the DOM, a trie of words (up to 29 children for the Turkish alphabet), B-trees in databases (hundreds of children).
      \end{itemize}
    \end{column}
  \end{columns}
\end{frame}
% ~~~~~~~~~~~~~~~~~~~~~~~~~~~~~~~~~~~~~~ A frm

% ~~~~~~~~~~~~~~~~~~~~~~~~~~~~~~~~~~~~~~ V
\hSection[Representation]{Representing Trees}
% ~~~~~~~~~~~~~~~~~~~~~~~~~~~~~~~~~~~~~~ A




% ~~~~~~~~~~~~~~~~~~~~~~~~~~~~~~~~~~~~~~ V frm
\begin{frame}[t,fragile] \frametitle{Linked nodes}
  \begin{columns}[T]
    \begin{column}{0.5\textwidth}
      \begin{block}{Binary tree}


% +++++++++++++++++++++++++++++++++++++++ V lst
%: -Lst.
\begin{lstlisting}[style=codesnippet, language=Java, basicstyle=\ttfamily\scriptsize]
class Node {
    int value;
    Node left, right;     // null: no child
}
\end{lstlisting}
% +++++++++++++++++++++++++++++++++++++++ A lst


      \end{block}
      \begin{block}{N-ary tree}


% +++++++++++++++++++++++++++++++++++++++ V lst
%: -Lst.
\begin{lstlisting}[style=codesnippet, language=Java, basicstyle=\ttfamily\scriptsize]
class Node {
    String name;
    List<Node> children = new ArrayList<>();
}
\end{lstlisting}
% +++++++++++++++++++++++++++++++++++++++ A lst


      \end{block}
    \end{column}
    \begin{column}{0.46\textwidth}
      \begin{center}
      \begin{tikzpicture}[box/.style={draw, thick, minimum width=0.55cm, minimum height=0.45cm, font=\small}]
        \node[box, fill=colorTop!50] (a) at (0,0) {A};
        \node[box] (al) at (-0.55,0) {}; \node[box] (ar) at (0.55,0) {};
        \node[box, fill=red!20] (b) at (-1.6,-1.4) {B};
        \node[box] (bl) at (-2.15,-1.4) {}; \node[box] (br) at (-1.05,-1.4) {};
        \node[box, fill=red!20] (c) at (1.6,-1.4) {C};
        \node[box] (cl) at (1.05,-1.4) {}; \node[box] (cr) at (2.15,-1.4) {};
        \draw[-latex, thick] (al.center) -- (b.north);
        \draw[-latex, thick] (ar.center) -- (c.north);
        \node[font=\tiny] at (-2.15,-1.9) {null}; \node[font=\tiny] at (-1.05,-1.9) {null};
        \node[font=\tiny] at (1.05,-1.9) {null};  \node[font=\tiny] at (2.15,-1.9) {null};
      \end{tikzpicture}
      \end{center}
      \small Each node: a value and two references. A missing child is \texttt{null}.
      The tree is just a reference to the root.
    \end{column}
  \end{columns}
\end{frame}
% ~~~~~~~~~~~~~~~~~~~~~~~~~~~~~~~~~~~~~~ A frm




% ~~~~~~~~~~~~~~~~~~~~~~~~~~~~~~~~~~~~~~ V frm
\begin{frame}[t,fragile] \frametitle{A complete binary tree fits in an array}
  \begin{columns}[T]
    \begin{column}{0.4\textwidth}
      \begin{center}
      \begin{tikzpicture}[htree, level distance=0.95cm, level 1/.style={sibling distance=2.0cm}, level 2/.style={sibling distance=1.0cm}, scale=1.0, transform shape]
  \node[hnode, horder=0] {A} child {node[hnode, horder=1] {B} child {node[hnode, horder=3] {D}} child {node[hnode, horder=4] {E}}} child {node[hnode, horder=2] {C} child {node[hnode, horder=5] {F}} child[missing]};
\end{tikzpicture}
      \end{center}
    \end{column}
    \begin{column}{0.56\textwidth}
      \begin{center}
      \begin{tikzpicture}\hString[0.55]{A,B,C,D,E,F}\end{tikzpicture}
      \end{center}
      Store the nodes level by level; for the node at index $i$:
      \[
      	\mathrm{left}(i) = 2i + 1, \qquad \mathrm{right}(i) = 2i + 2
      \]
      \[
      	\mathrm{parent}(i) = \left\lfloor \frac{i - 1}{2} \right\rfloor
      \]
      \begin{itemize}\small
        \item No references at all: compact and fast.
        \item Only for \textbf{complete} trees; gaps waste space. Used by heaps (\texttt{PriorityQueue}).
      \end{itemize}
    \end{column}
  \end{columns}
\end{frame}
% ~~~~~~~~~~~~~~~~~~~~~~~~~~~~~~~~~~~~~~ A frm




% ~~~~~~~~~~~~~~~~~~~~~~~~~~~~~~~~~~~~~~ V frm
\begin{frame}[t,fragile] \frametitle{N-ary trees as binary trees: first child, next sibling}
  \begin{columns}[T]
    \begin{column}{0.44\textwidth}
      \begin{center}
      \begin{tikzpicture}[htree, level distance=0.95cm, level 1/.style={sibling distance=2.0cm}, level 2/.style={sibling distance=0.75cm}, scale=1.0, transform shape]
  \node[hnode] {R} child {node[hnode] {A} child {node[hnode] {D}} child {node[hnode] {E}}} child {node[hnode] {B}} child {node[hnode] {C} child {node[hnode] {F}} child {node[hnode] {G}} child {node[hnode] {H}}};
\end{tikzpicture}
      \end{center}
    \end{column}
    \begin{column}{0.54\textwidth}
      \begin{center}
      \begin{tikzpicture}[every node/.style={hnode}, scale=0.85, transform shape]
        \node (r) at (0,0) {R};
        \node (a) at (0,-1) {A}; \node (b) at (1.2,-1) {B}; \node (c) at (2.4,-1) {C};
        \node (d) at (0,-2) {D}; \node (e) at (1.2,-2) {E};
        \node (f) at (2.4,-2) {F}; \node (g) at (3.6,-2) {G}; \node (h) at (4.8,-2) {H};
        \draw[-latex, thick, blue] (r) -- (a); \draw[-latex, thick, blue] (a) -- (d); \draw[-latex, thick, blue] (c) -- (f);
        \draw[-latex, thick, red] (a) -- (b); \draw[-latex, thick, red] (b) -- (c);
        \draw[-latex, thick, red] (d) -- (e); \draw[-latex, thick, red] (f) -- (g); \draw[-latex, thick, red] (g) -- (h);
      \end{tikzpicture}
      \end{center}
      \small \textcolor{blue}{blue}: first child, \textcolor{red}{red}: next sibling. Every node needs only \textbf{two} references,
      whatever its number of children: an N-ary tree stored as a binary tree.
    \end{column}
  \end{columns}
\end{frame}
% ~~~~~~~~~~~~~~~~~~~~~~~~~~~~~~~~~~~~~~ A frm

% ~~~~~~~~~~~~~~~~~~~~~~~~~~~~~~~~~~~~~~ V
\hSection[Traversal]{Tree Traversal}
% ~~~~~~~~~~~~~~~~~~~~~~~~~~~~~~~~~~~~~~ A

% ~~~~~~~~~~~~~~~~~~~~~~~~~~~~~~~~~~~~~~ V
\subsection[Depth first]{Depth first}
% ~~~~~~~~~~~~~~~~~~~~~~~~~~~~~~~~~~~~~~ A




% ~~~~~~~~~~~~~~~~~~~~~~~~~~~~~~~~~~~~~~ V frm
\begin{frame}[t,fragile] \frametitle{Visiting every node: three depth-first orders}
  \begin{columns}[T]
    \begin{column}{0.32\textwidth}
      \centering
      \begin{tikzpicture}[htree, level distance=0.95cm, level 1/.style={sibling distance=1.8cm}, level 2/.style={sibling distance=0.9cm}, level 3/.style={sibling distance=0.5cm}, scale=0.8, transform shape]
  \node[hnode, horder=1] {A} child {node[hnode, horder=2] {B} child {node[hnode, horder=3] {D}} child {node[hnode, horder=4] {E}}} child {node[hnode, horder=5] {C} child[missing] child {node[hnode, horder=6] {F} child {node[hnode, horder=7] {G}} child[missing]}};
\end{tikzpicture}\\
      \textbf{preorder}: node, left, right\\ \small\texttt{A B D E C F G}
    \end{column}
    \begin{column}{0.32\textwidth}
      \centering
      \begin{tikzpicture}[htree, level distance=0.95cm, level 1/.style={sibling distance=1.8cm}, level 2/.style={sibling distance=0.9cm}, level 3/.style={sibling distance=0.5cm}, scale=0.8, transform shape]
  \node[hnode, horder=4] {A} child {node[hnode, horder=2] {B} child {node[hnode, horder=1] {D}} child {node[hnode, horder=3] {E}}} child {node[hnode, horder=5] {C} child[missing] child {node[hnode, horder=7] {F} child {node[hnode, horder=6] {G}} child[missing]}};
\end{tikzpicture}\\
      \textbf{inorder}: left, node, right\\ \small\texttt{D B E A C G F}
    \end{column}
    \begin{column}{0.32\textwidth}
      \centering
      \begin{tikzpicture}[htree, level distance=0.95cm, level 1/.style={sibling distance=1.8cm}, level 2/.style={sibling distance=0.9cm}, level 3/.style={sibling distance=0.5cm}, scale=0.8, transform shape]
  \node[hnode, horder=7] {A} child {node[hnode, horder=3] {B} child {node[hnode, horder=1] {D}} child {node[hnode, horder=2] {E}}} child {node[hnode, horder=6] {C} child[missing] child {node[hnode, horder=5] {F} child {node[hnode, horder=4] {G}} child[missing]}};
\end{tikzpicture}\\
      \textbf{postorder}: left, right, node\\ \small\texttt{D E B G F C A}
    \end{column}
  \end{columns}
  \vspace{0.6em}
  \small The red numbers give the visit order. The name says \textbf{when the node itself} is visited: before (\emph{pre}), between (\emph{in}) or after (\emph{post}) its subtrees.
\end{frame}
% ~~~~~~~~~~~~~~~~~~~~~~~~~~~~~~~~~~~~~~ A frm




% ~~~~~~~~~~~~~~~~~~~~~~~~~~~~~~~~~~~~~~ V frm
\begin{frame}[t,fragile] \frametitle{The three orders as algorithms}
  \begin{columns}[T]
    \begin{column}{0.32\textwidth}
      \SetAlFnt{\small}


% +++++++++++++++++++++++++++++++++++++++ V alg
%: -Alg. 
\begin{algorithm}[H]
  \caption{Preorder($t$)}
  \lIf{$t$ is empty}{\Return}
  visit($t$)\;
  Preorder($t.left$)\;
  Preorder($t.right$)\;
\end{algorithm}
% +++++++++++++++++++++++++++++++++++++++ A alg

    \end{column}
    \begin{column}{0.32\textwidth}
      \SetAlFnt{\small}


% +++++++++++++++++++++++++++++++++++++++ V alg
%: -Alg. 
\begin{algorithm}[H]
  \caption{Inorder($t$)}
  \lIf{$t$ is empty}{\Return}
  Inorder($t.left$)\;
  visit($t$)\;
  Inorder($t.right$)\;
\end{algorithm}
% +++++++++++++++++++++++++++++++++++++++ A alg

    \end{column}
    \begin{column}{0.32\textwidth}
      \SetAlFnt{\small}


% +++++++++++++++++++++++++++++++++++++++ V alg
%: -Alg. 
\begin{algorithm}[H]
  \caption{Postorder($t$)}
  \lIf{$t$ is empty}{\Return}
  Postorder($t.left$)\;
  Postorder($t.right$)\;
  visit($t$)\;
\end{algorithm}
% +++++++++++++++++++++++++++++++++++++++ A alg

    \end{column}
  \end{columns}
  \begin{itemize}\small
    \item Only the position of \texttt{visit} changes. Each node is visited once: $O(n)$ time.
    \item The recursion uses the call stack: $O(h)$ extra memory, $h$ the height.
  \end{itemize}
\end{frame}
% ~~~~~~~~~~~~~~~~~~~~~~~~~~~~~~~~~~~~~~ A frm




% ~~~~~~~~~~~~~~~~~~~~~~~~~~~~~~~~~~~~~~ V frm
\begin{frame}[t,fragile] \frametitle{A trick: walk around the tree}
  \begin{columns}[T]
    \begin{column}{0.46\textwidth}
      \begin{center}
      \begin{tikzpicture}[htree, level distance=1.2cm, level 1/.style={sibling distance=2.6cm}, level 2/.style={sibling distance=1.3cm},
                          dot/.style={circle, inner sep=1.6pt}]
        \node[hnode] (A) {A}
          child {node[hnode] (B) {B} child {node[hnode] (D) {D}} child {node[hnode] (E) {E}}}
          child {node[hnode] (C) {C} child[missing] child {node[hnode] (F) {F}}};
        \foreach \n in {A, B, C, D, E, F} {
          \node[dot, fill=blue]  at ($(\n)+(-0.45,0)$) {};
          \node[dot, fill=green!60!black] at ($(\n)+(0,-0.45)$) {};
          \node[dot, fill=red]   at ($(\n)+(0.45,0)$) {};
        }
      \end{tikzpicture}
      \end{center}
    \end{column}
    \begin{column}{0.5\textwidth}
      Walk around the outside of the tree, counter-clockwise, starting left of the root (the \textbf{Euler tour}). You pass each node three times:
      \begin{itemize}\small
        \item its \textcolor{blue}{\textbf{left}} dot: preorder \texttt{A B D E C F}
        \item its \textcolor{green!60!black}{\textbf{bottom}} dot: inorder \texttt{D B E A C F}
        \item its \textcolor{red}{\textbf{right}} dot: postorder \texttt{D E B F C A}
      \end{itemize}
      \small Useful for checking an answer by hand.
    \end{column}
  \end{columns}
\end{frame}
% ~~~~~~~~~~~~~~~~~~~~~~~~~~~~~~~~~~~~~~ A frm

% ~~~~~~~~~~~~~~~~~~~~~~~~~~~~~~~~~~~~~~ V
\subsection[Breadth first]{Breadth first}
% ~~~~~~~~~~~~~~~~~~~~~~~~~~~~~~~~~~~~~~ A




% ~~~~~~~~~~~~~~~~~~~~~~~~~~~~~~~~~~~~~~ V frm
\begin{frame}[t,fragile] \frametitle{Level order: breadth first, with a queue}
  \begin{columns}[T]
    \begin{column}{0.36\textwidth}
      \centering
      \begin{tikzpicture}[htree, level distance=0.95cm, level 1/.style={sibling distance=1.8cm}, level 2/.style={sibling distance=0.9cm}, level 3/.style={sibling distance=0.5cm}, scale=0.85, transform shape]
  \node[hnode, horder=1] {A} child {node[hnode, horder=2] {B} child {node[hnode, horder=4] {D}} child {node[hnode, horder=5] {E}}} child {node[hnode, horder=3] {C} child[missing] child {node[hnode, horder=6] {F} child {node[hnode, horder=7] {G}} child[missing]}};
\end{tikzpicture}\\
      \small level order: \texttt{A B C D E F G}
    \end{column}
    \begin{column}{0.62\textwidth}
      \SetAlFnt{\small}


% +++++++++++++++++++++++++++++++++++++++ V alg
%: -Alg. 
\begin{algorithm}[H]
  \caption{LevelOrder($t$)}
  \lIf{$t$ is empty}{\Return}
  $Q \gets$ a queue containing $t$\;
  \While{$Q$ is not empty}{
    $n \gets$ \textsc{Dequeue}$(Q)$\;
    visit($n$)\;
    \lForEach{child $c$ of $n$, from left to right}{\textsc{Enqueue}$(Q, c)$}
  }
\end{algorithm}
% +++++++++++++++++++++++++++++++++++++++ A alg

      \small A \textbf{queue} (FIFO) gives level order; the same loop with a \textbf{stack} (LIFO) gives a depth-first order.
    \end{column}
  \end{columns}
\end{frame}
% ~~~~~~~~~~~~~~~~~~~~~~~~~~~~~~~~~~~~~~ A frm




% ~~~~~~~~~~~~~~~~~~~~~~~~~~~~~~~~~~~~~~ V frm
\begin{frame}[t,fragile] \frametitle{Traversing N-ary trees}
  \begin{columns}[T]
    \begin{column}{0.4\textwidth}
      \begin{center}
      \begin{tikzpicture}[htree, level distance=0.95cm, level 1/.style={sibling distance=1.8cm}, level 2/.style={sibling distance=0.7cm}, scale=1.0, transform shape]
  \node[hnode, horder=1] {R} child {node[hnode, horder=2] {A} child {node[hnode, horder=3] {D}} child {node[hnode, horder=4] {E}}} child {node[hnode, horder=5] {B}} child {node[hnode, horder=6] {C} child {node[hnode, horder=7] {F}} child {node[hnode, horder=8] {G}} child {node[hnode, horder=9] {H}}};
\end{tikzpicture}
      \end{center}
      \small preorder: \texttt{R A D E B C F G H}\\
      postorder: \texttt{D E A B F G H C R}\\
      level order: \texttt{R A B C D E F G H}
    \end{column}
    \begin{column}{0.56\textwidth}
      \SetAlFnt{\small}


% +++++++++++++++++++++++++++++++++++++++ V alg
%: -Alg. 
\begin{algorithm}[H]
  \caption{Preorder and postorder of an N-ary tree}
  \Fn{\textsc{Pre}($n$)}{
    visit($n$)\;
    \lForEach{child $c$ of $n$}{\textsc{Pre}($c$)}
  }
  \Fn{\textsc{Post}($n$)}{
    \lForEach{child $c$ of $n$}{\textsc{Post}($c$)}
    visit($n$)\;
  }
\end{algorithm}
% +++++++++++++++++++++++++++++++++++++++ A alg

      \small There is no natural \emph{inorder}: with three children, ``between'' is not defined.
    \end{column}
  \end{columns}
\end{frame}
% ~~~~~~~~~~~~~~~~~~~~~~~~~~~~~~~~~~~~~~ A frm

% ~~~~~~~~~~~~~~~~~~~~~~~~~~~~~~~~~~~~~~ V
\hSection[Uses]{Which Traversal for What?}
% ~~~~~~~~~~~~~~~~~~~~~~~~~~~~~~~~~~~~~~ A




% ~~~~~~~~~~~~~~~~~~~~~~~~~~~~~~~~~~~~~~ V frm
\begin{frame}[t,fragile] \frametitle{An expression tree, three ways}
  \begin{columns}[T]
    \begin{column}{0.36\textwidth}
      \begin{center}
      \begin{tikzpicture}[htree, level distance=0.95cm, level 1/.style={sibling distance=2.0cm}, level 2/.style={sibling distance=1.0cm}, scale=1.0, transform shape]
  \node[hsquare] {*} child {node[hsquare] {+} child {node[hnode] {3}} child {node[hnode] {4}}} child {node[hsquare] {-} child {node[hnode] {5}} child {node[hnode] {2}}};
\end{tikzpicture}
      \end{center}
    \end{column}
    \begin{column}{0.6\textwidth}
      \begin{center}\small
      \begin{tabular}{lll}
        \toprule
        \textbf{traversal} & \textbf{gives}            & \textbf{notation} \\
        \midrule
        preorder           & \texttt{* + 3 4 - 5 2}    & prefix (Polish) \\
        inorder            & \texttt{3 + 4 * 5 - 2}    & infix, needs brackets \\
        postorder          & \texttt{3 4 + 5 2 - *}    & postfix (reverse Polish) \\
        \bottomrule
      \end{tabular}
      \end{center}
      \begin{itemize}\small
        \item To \textbf{evaluate}, a node needs the values of its children first: postorder. $(3 + 4) \times (5 - 2) = 21$.
        \item Inorder must add brackets around each operator node: \texttt{((3 + 4) * (5 - 2))}.
        \item Compilers build such trees (\emph{abstract syntax trees}) from your Java code.
      \end{itemize}
    \end{column}
  \end{columns}
\end{frame}
% ~~~~~~~~~~~~~~~~~~~~~~~~~~~~~~~~~~~~~~ A frm




% ~~~~~~~~~~~~~~~~~~~~~~~~~~~~~~~~~~~~~~ V frm
\begin{frame}[t,fragile] \frametitle{Choosing a traversal}
  \begin{center}\small
  \begin{tabular}{lL{0.5\textwidth}}
    \toprule
    \textbf{traversal} & \textbf{when the task \dots} \\
    \midrule
    preorder           & handles a node \textbf{before} its children: print an outline with indentation, copy a tree, \texttt{mkdir} folders \\
    inorder            & needs a binary search tree's keys \textbf{in sorted order} \\
    postorder          & needs the children's results \textbf{first}: size of a folder, height, evaluate, delete a tree \\
    level order        & works \textbf{level by level}: nearest first, print by levels, shortest number of steps \\
    \bottomrule
  \end{tabular}
  \end{center}
  \vspace{0.3em}
  \begin{center}\small
  \begin{tabular}{lll}
    \toprule
                 & \textbf{time} & \textbf{extra memory} \\
    \midrule
    depth first (recursion or a stack) & $O(n)$ & $O(h)$: the height \\
    breadth first (a queue)            & $O(n)$ & $O(w)$: the widest level \\
    \bottomrule
  \end{tabular}
  \end{center}
\end{frame}
% ~~~~~~~~~~~~~~~~~~~~~~~~~~~~~~~~~~~~~~ A frm

% ~~~~~~~~~~~~~~~~~~~~~~~~~~~~~~~~~~~~~~ V
\hSection[Java]{Trees in Java}
% ~~~~~~~~~~~~~~~~~~~~~~~~~~~~~~~~~~~~~~ A




% ~~~~~~~~~~~~~~~~~~~~~~~~~~~~~~~~~~~~~~ V frm
\begin{frame}[t,fragile] \frametitle{Where Java uses trees}
  \begin{center}\small
  \begin{tabular}{lL{0.5\textwidth}}
    \toprule
    \textbf{class}                        & \textbf{tree} \\
    \midrule
    \texttt{java.util.TreeMap}, \texttt{TreeSet} & a red-black tree: a balanced binary search tree, keys kept sorted \\
    \texttt{java.util.PriorityQueue}      & a binary heap: a complete binary tree in an array \\
    \texttt{java.nio.file.Files.walk}     & visits the file tree (preorder) \\
    \texttt{org.w3c.dom.Node}             & an XML document as an N-ary tree \\
    \texttt{javax.swing.JTree}            & shows a tree in a window \\
    \bottomrule
  \end{tabular}
  \end{center}
  \vspace{0.3em}
  \small There is no general \texttt{Tree} class in the JDK: for your own data you write a small \texttt{Node} class (next decks).
\end{frame}
% ~~~~~~~~~~~~~~~~~~~~~~~~~~~~~~~~~~~~~~ A frm




% ~~~~~~~~~~~~~~~~~~~~~~~~~~~~~~~~~~~~~~ V frm
\begin{frame}[t,fragile] \frametitle{Trying it: a \texttt{TreeMap} in \texttt{jshell}}
  \small Open a terminal (\win: \emph{Command Prompt}; \mac: \emph{Terminal}), type \texttt{jshell}, then these lines (\texttt{try.jsh}):
  \lstcode[basicstyle=\ttfamily\tiny, numbers=none]{codeoutput/JShell/try.jsh}{Java}

  \lstterminal[basicstyle=\ttfamily\footnotesize]{codeoutput/JShell.txt}

  \vspace{0.1em}
  \small
  The keys come out \textbf{sorted}: \texttt{TreeMap} walks its tree in inorder.
\end{frame}
% ~~~~~~~~~~~~~~~~~~~~~~~~~~~~~~~~~~~~~~ A frm

% ~~~~~~~~~~~~~~~~~~~~~~~~~~~~~~~~~~~~~~ V
\hSection[Summary]{Summary}
% ~~~~~~~~~~~~~~~~~~~~~~~~~~~~~~~~~~~~~~ A




% ~~~~~~~~~~~~~~~~~~~~~~~~~~~~~~~~~~~~~~ V frm
\begin{frame}[t,fragile] \frametitle{Big ideas}
  \begin{itemize}
    \item A tree is a hierarchy: one root, every other node has exactly one parent, no cycles; $n$ nodes, $n - 1$ edges.
    \item Binary trees: left and right child. Height between $\lfloor \log_{2} n \rfloor$ and $n - 1$.
    \item N-ary trees: a list of children, or first child / next sibling.
    \item Complete binary trees fit in an array: children at $2i + 1$, $2i + 2$.
    \item Traversals: preorder, inorder, postorder (depth first, $O(h)$ memory), level order (queue, $O(w)$ memory); all $O(n)$.
    \item Pick the traversal by \textbf{when} a node must be handled relative to its children.
  \end{itemize}
\end{frame}
% ~~~~~~~~~~~~~~~~~~~~~~~~~~~~~~~~~~~~~~ A frm




% ~~~~~~~~~~~~~~~~~~~~~~~~~~~~~~~~~~~~~~ V frm
\begin{frame}[t,fragile] \frametitle{Check yourself}
  \begin{columns}[T]
    \begin{column}{0.36\textwidth}
      \begin{center}
      \begin{tikzpicture}[htree, level distance=0.95cm, level 1/.style={sibling distance=2.0cm}, level 2/.style={sibling distance=1.0cm}, level 3/.style={sibling distance=0.6cm}, scale=1.0, transform shape]
  \node[hnode] {M} child {node[hnode] {F} child {node[hnode] {B}} child {node[hnode] {H} child {node[hnode] {G}} child[missing]}} child {node[hnode] {S} child[missing] child {node[hnode] {W}}};
\end{tikzpicture}
      \end{center}
    \end{column}
    \begin{column}{0.6\textwidth}
      \begin{enumerate}\small
        \item Give the root, the leaves, the height, and the depth of \texttt{G}.
        \item Write the preorder, inorder, postorder and level order.
        \item Is it full? complete? Why is its inorder sorted?
        \item Store it in an array as if it were complete. Which slots stay empty?
        \item Draw the expression tree of \texttt{2 * (3 + 4) - 10 / 5}; write its postfix form.
      \end{enumerate}
    \end{column}
  \end{columns}
\end{frame}
% ~~~~~~~~~~~~~~~~~~~~~~~~~~~~~~~~~~~~~~ A frm




% ~~~~~~~~~~~~~~~~~~~~~~~~~~~~~~~~~~~~~~ V frm
\begin{frame}[t,fragile]
	\begin{center}
		\vspace{4cm}
		\hRemark{\Huge Questions?}
	\end{center}
\end{frame}
% ~~~~~~~~~~~~~~~~~~~~~~~~~~~~~~~~~~~~~~ A frm



\end{document}
